% Source: https://artofproblemsolving.com/community/c3h2854244_geometry_problem
% Stable identity: aops-topic-2854244-q01
Since $AB$ is a diameter, $\angle ACB=90^\circ$, so the bisector $CD$ makes $\angle BCD=45^\circ$. Also $\angle CBD=90^\circ$. So triangle $CBD$ is an isosceles right triangle with hypotenuse $CD$ and $BC=BD$. The median $BM$ to the hypotenuse of an isosceles triangle is also an altitude, so $\angle BMD=90^\circ$.
\begin{asy}
usepackage("amsmath");
size(5cm);
defaultpen(fontsize(10pt));
pair A=(-3,0),B=(3,0),C=3*dir(115);
draw(circle((0,0),3),gray);draw(A--B--C--cycle);
pair u=unit(unit(A-C)+unit(B-C));pair Bp=B+rotate(90)*(B-C);
pair D=extension(C,C+u,B,Bp);pair E=A+dot(B-A,unit(D-A))*unit(D-A);pair M=(C+D)/2;
draw(C--D);draw(B--D);draw(A--D,dashed);draw(B--E,dashed);draw(B--M,dotted);
draw(circle((B+D)/2,abs(B-D)/2),blue);
dot(E);dot(M);label("$A$",A,W);label("$B$",B,E);label("$C$",C,N);label("$D$",D,S);label("$E$",E,SW);label("$M$",M,W);
\end{asy}

By construction $\angle BED=90^\circ$. So $M$ and $E$ both see the segment $BD$ at a right angle, and by the converse of Thales' theorem they lie on the circle with diameter $BD$. Hence $B,D,E,M$ are concyclic. $\blacksquare$
